% ch5.tex -- Module 8 ESSENTIALS edition, Chapter 5: Deep inside the atom (IQ5).
% book): same notation, same facts and numbers. All figures are TikZ copied (and shrunk)
% from the full chapter, so no figure files are needed.
\chapter[Deep Inside the Atom]{Deep Inside the Atom: the Standard Model and Particle Accelerators}\label{ch5:chap}
\covers{M8.IQ5.01, M8.IQ5.02, M8.IQ5.03}

% chapter-local macros (same as the full book)
\newcommand\chFiveaq[1]{\ensuremath{\overline{\mathrm{#1}}}} % antiquark, e.g. \chFiveaq{d}
\newcommand\chFiveperc{\ensuremath{/c^2}}                     % "GeV/c^2" after \qty{..}{GeV}
\newcommand\chFiveanu{\ensuremath{\overline{\nu}_\mathrm{e}}} % electron antineutrino

\begin{summarybox}{In this chapter}
\begin{itemize}
  \item Evidence that protons and neutrons contain quarks; other particles.
  \item The Standard Model: quarks, leptons, hadrons and the four forces.
  \item How accelerators work ($W = qV$, $r = mv/(qB)$) and test the Standard Model.
\end{itemize}
\end{summarybox}

%% ==================================================================
\section{Protons and neutrons are not fundamental}\label{ch5:sec:evidence}
\covers{M8.IQ5.01}

\begin{bigidea}
A \term{fundamental particle} has no smaller parts. The electron is fundamental, but
protons and neutrons are made of quarks.
\end{bigidea}

Three pieces of evidence:
\begin{itemize}
  \item \textbf{Deep inelastic scattering} (SLAC, \cref{ch5:sec:experiments}): fast
        electrons broke protons apart. Far more electrons bounced at large angles than a
        smooth proton allows (\cref{ch5:fig:dis}), so the proton has tiny, \term{point-like}
        charges (like dots): \term{quarks}.
\end{itemize}

\begin{figure}
  \centering
  \begin{tikzpicture}[x=0.72cm, y=0.72cm,
      etrack/.style={line width=0.8pt, color=hscvel, midarrow},
      eout/.style={line width=0.8pt, color=hscvel, -{Stealth[length=2mm, width=1.5mm]}},
      qk/.style={circle, text=white, font=\scriptsize\bfseries, inner sep=0pt,
        minimum size=0.36cm}]
    \node[font=\footnotesize\bfseries, text=hscink] at (0,1.65) {Model A: smooth proton};
    \fill[hscforce!14] (0,0) circle (1.05);
    \draw[hscforce!70, line width=0.7pt] (0,0) circle (1.05);
    \draw[etrack] (-2.4,0.75) -- (-0.4,0.75) -- (2.4,0.55);
    \draw[etrack] (-2.4,0.35) -- (-0.2,0.35) -- (2.4,0.24);
    \draw[etrack] (-2.4,-0.45) -- (-0.2,-0.45) -- (2.4,-0.31);
    \draw[etrack] (-2.4,-0.85) -- (-0.4,-0.85) -- (2.4,-0.64);
    \node[font=\scriptsize, text=hscforce!80!black] at (0,-1.38)
      {$+$ charge spread evenly};
    \node[font=\scriptsize, text=hscmuted] at (0,-1.85) {prediction: only small deflections};
    \node[font=\scriptsize, text=hscvel, anchor=east] at (-2.45,0) {electrons};
    \begin{scope}[xshift=5.2cm]
      \node[font=\footnotesize\bfseries, text=hscink] at (0,1.65) {Model B: tiny charges inside};
      \draw[hscmuted, dashed, line width=0.7pt] (0,0) circle (1.05);
      \draw[etrack] (-2.4,0.20) -- (0,0.20) -- (2.4,0.13);
      \draw[etrack] (-2.4,0.70) -- (0.07,0.70);
      \draw[eout] (0.07,0.70) -- (-0.97,1.30);
      \draw[etrack] (-2.4,-0.62) -- (0.10,-0.62);
      \draw[eout] (0.10,-0.62) -- (0.60,-1.49);
      \node[qk, fill=hscseriesC!85!black] at (0.30,0.70) {d};
      \node[qk, fill=hscseriesB] at (-0.45,-0.20) {u};
      \node[qk, fill=hscseriesB] at (0.35,-0.62) {u};
      \node[font=\scriptsize, text=hscforce!80!black] at (0,-1.85)
        {SLAC: far more scatter at large angles};
    \end{scope}
  \end{tikzpicture}
  \caption{Two models of the proton. Large-angle bounces need point-like charges inside:
    quarks (u = up, d = down).}\label{ch5:fig:dis}
\end{figure}

\begin{itemize}
  \item \textbf{Beta decay:} $\mathrm{n} \to \mathrm{p} + \mathrm{e}^- + \chFiveanu$
        turns udd into uud: just \emph{one} down quark changes into an up quark. This is easy
        to explain if the neutron has parts, so it \emph{supports} the quark model. It is not
        proof on its own: the muon is fundamental, but it also decays.
  \item \textbf{The particle zoo:} hundreds of new particles (1950s--60s). In 1964
        Gell-Mann and Zweig built them all from three quarks: up, down, strange. Gell-Mann's
        pattern (1962) predicted the $\Omega^-$ (sss), found in 1964 as predicted.
\end{itemize}

\subsection{Other subatomic particles}

Each particle in \cref{ch5:tab:particles} was new: no known particle had its mass, charge
or behaviour. An \term{antiparticle} has its partner's mass but opposite charge.

\begin{table}
  \caption{Some particles that are not protons, neutrons or electrons.}\label{ch5:tab:particles}
  \centering
  \small
  \begin{tabular}{@{}L{42mm}L{97mm}@{}}
    \toprule
    \colhead{Particle} & \colhead{What was seen, and what it shows}\\
    \midrule
    positron $\mathrm{e}^+$ (1932): the first antiparticle & Anderson's \term{cloud chamber}
      (charged particles leave trails of drops) in a magnetic field: a track like an
      electron's, but curving the other way (a lead plate showed its direction), so
      positive. Predicted by Dirac (1928).\\
    muon $\mu^-$ (1936): a heavy electron & Tracks of \term{cosmic rays} (fast particles
      from space): charge $-e$, about 207 electron masses.\\
    pions $\pi^\pm$ (1947): mesons & Cosmic-ray tracks: about 270 electron masses.\\
    antiproton $\overline{\mathrm{p}}$ (1955) & Made in an accelerator: proton's mass,
      charge $-e$.\\
    neutrino $\nu$ (1956): neutral, tiny mass & Beta electrons have a spread of energies:
      Pauli (1930) said a hidden particle carries the missing energy. Detected near a
      nuclear reactor.\\
    \bottomrule
  \end{tabular}
\end{table}

%% ==================================================================
\section{The Standard Model: particles and forces}\label{ch5:sec:sm}
\covers{M8.IQ5.02}

\begin{bigidea}
The \term{Standard Model} is the theory of the fundamental particles and forces
(\cref{ch5:fig:sm}): 6 quarks, 6 leptons and the bosons that carry the forces.
\end{bigidea}

\begin{keywords}
  \kw{Quark}{a fundamental matter particle (fermion) that feels the strong force.}
  \kw{Lepton}{a fundamental matter particle that does not feel the strong force.}
  \kw{Hadron}{a particle made of quarks, such as a proton.}
  \kw{Fermion / boson}{half-number / whole-number spin (built-in turning).}
\end{keywords}

\begin{itemize}
  \item Six quark \term{flavours} (types) in three \term{generations}, each heavier than the
        one before; ordinary matter uses only u and d. \term{Colour charge}: the strong force's
        ``charge'' (not real colour).
  \item \term{Confinement}: a quark is never found alone, only inside hadrons. Antiquarks
        have a bar: $\chFiveaq{u}$ is $-\tfrac23 e$, $\chFiveaq{d}$ is $+\tfrac13 e$.
  \item Neutrinos have a tiny mass and no electric or colour charge: only the weak force and
        gravity act on them.
\end{itemize}

\begin{figure}
  \centering
  \begin{tikzpicture}[x=0.86cm, y=0.8cm,
      cell/.style={draw=#1!75!black, fill=#1!14, rounded corners=1.5mm,
        minimum width=1.62cm, minimum height=1.08cm, line width=0.6pt,
        align=center, inner sep=0.5pt, font=\scriptsize,
        execute at begin node={\setstretch{0.85}}},
      bcell/.style={cell=#1, minimum width=1.9cm}]
    \foreach \x/\g in {0/1, 2.05/2, 4.1/3}
      \node[font=\scriptsize\bfseries] at (\x,1.08) {Generation \g};
    \node[font=\scriptsize\bfseries] at (6.55,1.08) {Force carriers};
    \node[font=\scriptsize\bfseries] at (8.95,1.08) {Higgs};
    \node[cell=hscseriesB] at (0,0)    {{\small $\mathrm{u}$}\\up\\$+\tfrac23$};
    \node[cell=hscseriesB] at (2.05,0) {{\small $\mathrm{c}$}\\charm\\$+\tfrac23$};
    \node[cell=hscseriesB] at (4.1,0)  {{\small $\mathrm{t}$}\\top\\$+\tfrac23$};
    \node[cell=hscseriesB] at (0,-1.48)    {{\small $\mathrm{d}$}\\down\\$-\tfrac13$};
    \node[cell=hscseriesB] at (2.05,-1.48) {{\small $\mathrm{s}$}\\strange\\$-\tfrac13$};
    \node[cell=hscseriesB] at (4.1,-1.48)  {{\small $\mathrm{b}$}\\bottom\\$-\tfrac13$};
    \node[cell=hscseriesC] at (0,-2.96)    {{\small $\mathrm{e}^-$}\\electron\\$-1$};
    \node[cell=hscseriesC] at (2.05,-2.96) {{\small $\mu^-$}\\muon\\$-1$};
    \node[cell=hscseriesC] at (4.1,-2.96)  {{\small $\tau^-$}\\tau\\$-1$};
    \node[cell=hscseriesC] at (0,-4.44)    {{\small $\nu_\mathrm{e}$}\\e-neutrino\\$0$};
    \node[cell=hscseriesC] at (2.05,-4.44) {{\small $\nu_\mu$}\\$\mu$-neutrino\\$0$};
    \node[cell=hscseriesC] at (4.1,-4.44)  {{\small $\nu_\tau$}\\$\tau$-neutrino\\$0$};
    \node[bcell=hscseriesA] at (6.55,0)     {{\small $\mathrm{g}$}\\gluon\\strong force};
    \node[bcell=hscseriesA] at (6.55,-1.48) {{\small $\gamma$}\\photon\\EM force};
    \node[bcell=hscseriesA] at (6.55,-2.96) {{\small $\mathrm{W}^\pm$}\\W bosons\\weak force};
    \node[bcell=hscseriesA] at (6.55,-4.44) {{\small $\mathrm{Z}^0$}\\Z boson\\weak force};
    \node[cell=hscbrand] at (8.95,0) {{\small $\mathrm{H}$}\\Higgs boson\\spin 0};
    \node[draw=hscmuted, dashed, rounded corners=1.5mm, minimum width=1.62cm,
      minimum height=1.08cm, align=center, inner sep=0.5pt, line width=0.6pt,
      font=\scriptsize, execute at begin node={\setstretch{0.85}}] (grav) at (8.95,-2.35)
      {{\small\color{hscmuted} G?}\\{\color{hscmuted} graviton}};
    \node[font=\scriptsize, text=hscmuted, align=center, below=1pt of grav]
      {never found;\\ NOT in the\\ Standard Model};
    \draw[hscseriesB!75!black, line width=0.8pt] (-1.15,0.68) -- (-1.25,0.68) -- (-1.25,-2.16) -- (-1.15,-2.16);
    \node[rotate=90, font=\scriptsize\bfseries, text=hscseriesB!75!black] at (-1.55,-0.74) {QUARKS};
    \draw[hscseriesC!75!black, line width=0.8pt] (-1.15,-2.28) -- (-1.25,-2.28) -- (-1.25,-5.12) -- (-1.15,-5.12);
    \node[rotate=90, font=\scriptsize\bfseries, text=hscseriesC!75!black] at (-1.55,-3.70) {LEPTONS};
    \draw[hscmuted, line width=0.6pt] (-0.95,-5.3) -- (5.05,-5.3);
    \node[font=\scriptsize] at (2.05,-5.72) {matter particles: fermions (spin $\tfrac12$)};
    \draw[hscmuted, line width=0.6pt] (5.42,-5.3) -- (10.1,-5.3);
    \node[font=\scriptsize] at (7.75,-5.72) {bosons (whole-number spin)};
  \end{tikzpicture}
  \caption{The particles of the Standard Model. Numbers are charges in units of $e$.
    EM = electromagnetic.}\label{ch5:fig:sm}
\end{figure}

\subsection{The four fundamental forces}

An \term{exchange particle} is a boson that carries a force (\cref{ch5:tab:forces}), like a
ball thrown between two skaters that pushes them apart.

\begin{table}
  \caption{The four forces. Strengths are rough values compared with the strong force
    (books differ a little).}\label{ch5:tab:forces}
  \centering
  \small
  \begin{tabular}{@{}L{25mm}L{33mm}L{44mm}L{14mm}L{16mm}@{}}
    \toprule
    \colhead{Force} & \colhead{Carrier} & \colhead{Acts on} & \colhead{Strength}
      & \colhead{Range}\\
    \midrule
    strong & gluons (g) & quarks and gluons & 1 & \qty{e-15}{m}\\
    electromagnetic & photon ($\gamma$) & charged particles & $10^{-2}$ & no limit\\
    weak & $\mathrm{W}^+$, $\mathrm{W}^-$, $\mathrm{Z}^0$ & all quarks and leptons
      & $10^{-6}$ & \qty{e-18}{m}\\
    gravity & graviton? (not found) & all mass and energy & $10^{-38}$ & no limit\\
    \bottomrule
  \end{tabular}
\end{table}

\begin{itemize}
  \item \textbf{Strong:} a little of the pull between quarks leaks out of protons and
        neutrons. It holds the nucleus together, even though protons repel.
  \item \textbf{Weak:} the \emph{only} force that changes a quark's flavour. W and Z are very
        heavy (about 86 and 97 proton masses), so they live very briefly: tiny range.
  \item \textbf{Beta-minus:} $\mathrm{d} \to \mathrm{u} + \mathrm{W}^-$, then
        $\mathrm{W}^- \to \mathrm{e}^- + \chFiveanu$ (charge: $-\tfrac13 = +\tfrac23 - 1$).
        \textbf{Beta-plus} (only inside a nucleus): $\mathrm{u} \to \mathrm{d} +
        \mathrm{W}^+$, then $\mathrm{W}^+ \to \mathrm{e}^+ + \nu_\mathrm{e}$.
  \item \textbf{The Higgs field} gives most fundamental particles mass; its ripple is the
        \term{Higgs boson}.
\end{itemize}

\subsection{Hadrons: baryons and mesons}

\begin{itemize}
  \item A \term{baryon} is three quarks: proton (uud), neutron (udd). Three half-number
        spins add to a half-number, so baryons are fermions.
  \item A \term{meson} is a quark and an antiquark: pion $\pi^+$ (u$\chFiveaq{d}$). Two
        half-number spins add to a whole number, so mesons are bosons. All are unstable.
  \item \textbf{Charge of a hadron} = the sum of its quark charges (charge is conserved):
        $Q = q_1 + q_2 + q_3$ (baryon), $Q = q_1 + q_2$ (meson). \hscbadgenotonsheet\
        Flip the sign for an antiquark. The total is always a whole number of $e$.
\end{itemize}

%% ==================================================================
\section{Particle accelerators: how they work and what they show}\label{ch5:sec:accel}
\covers{M8.IQ5.03}

\begin{bigidea}
Electric fields speed charged particles up; magnetic fields steer them. Collisions of
fast particles test the Standard Model.
\end{bigidea}

\begin{itemize}
  \item Two reasons to accelerate: (1) a wave only shows details bigger than its wavelength; more
        momentum gives a shorter $\lambda = h/(mv)$ (near $c$, use $p_\mathrm{v}$). SLAC's
        $\lambda \approx \qty{6e-17}{m}$ was 27 times smaller than a proton. (2) Kinetic
        energy can become new mass ($E = mc^2$): a particle--antiparticle pair needs at
        least $2mc^2$.
  \item An electric field speeds a charge up: each gap adds $W = qV$ (on the sheet); after
        $n$ gaps from rest, $K = nqV$. $\qty{1}{eV} = \qty{1.602e-19}{J}$.
  \item The magnetic force $F = qv_\perp B = qvB\sin\theta$ is at \ang{90} to $v$: it turns
        the particle but does no work. Only \emph{charged} particles can be accelerated.
\end{itemize}

\subsection{The linear accelerator (linac)}

A \term{linac} (linear accelerator) is a straight row of hollow metal \term{drift tubes},
joined in turn to the two ends of an AC supply (\cref{ch5:fig:linac}).

\begin{figure}
  \centering
  \begin{tikzpicture}[x=0.85cm, y=0.85cm,
      tube/.style={draw=hscink, fill=hscmuted!18, line width=0.7pt}]
    \draw[draw=hscink, fill=hscbrand!10, line width=0.7pt] (-1.85,-0.38) rectangle (-0.5,0.38);
    \node[font=\scriptsize, align=center] at (-1.175,0) {proton\\[-2pt]source};
    \foreach \a/\b/\s in {0/0.850/+, 1.200/2.402/-, 2.752/4.224/+, 4.574/6.274/-, 6.624/8.525/+}{
      \draw[tube] (\a,-0.3) rectangle (\b,0.3);
      \node[font=\small\bfseries] at ({(\a+\b)/2},0) {$\s$};
    }
    \draw[guide] (-0.5,0) -- (0,0);
    \draw[velocity] (8.6,0) -- (9.6,0) node[midway, below, font=\scriptsize, text=hscvel] {beam};
    \fill[hscvel] (1.02,0) circle (2.2pt);
    \draw[force] (1.0,0.55) -- (1.6,0.55);
    \node[font=\scriptsize, text=hscforce, anchor=south] at (1.3,0.62) {push};
    \foreach \c in {0.425, 3.488, 7.575} \draw[hscink, line width=0.6pt] (\c,0.3) -- (\c,1.15);
    \foreach \c in {1.801, 5.424} \draw[hscink, line width=0.6pt] (\c,-0.3) -- (\c,-1.15);
    \draw[hscink, line width=0.6pt] (0.425,1.15) -- (7.575,1.15);
    \draw[hscink, line width=0.6pt] (4.6,1.15) -- (4.6,1.45);
    \draw[hscink, fill=white, line width=0.7pt] (4.6,1.75) circle (0.3);
    \node[font=\small] at (4.6,1.75) {$\sim$};
    \node[font=\scriptsize, anchor=west] at (4.95,1.75) {AC supply};
    \draw[hscink, line width=0.6pt] (4.6,2.05) -- (4.6,2.3) -- (-2.2,2.3) -- (-2.2,-1.15)
      -- (5.424,-1.15);
    \node[font=\scriptsize, anchor=north] at (3.6,-0.38) {no field in a tube};
    \draw[hscmuted, line width=0.5pt] (6.45,0.62) -- (6.45,0.36);
    \node[font=\scriptsize, anchor=south] at (6.15,0.6) {gap: field here};
    \draw[displacement] (2.75,-1.55) -- (8.5,-1.55) node[midway, below, font=\scriptsize,
      text=hscmuted] {each tube is longer than the one before};
  \end{tikzpicture}
  \caption{A linac at one instant. Half a cycle later every $+$ and $-$ swaps, so the
    next gap also pushes the proton forward.}\label{ch5:fig:linac}
\end{figure}

\begin{itemize}
  \item No field inside a tube: the particle coasts. Each gap pushes it forward; the AC
        swaps $+$ and $-$ in time for the next gap.
  \item It speeds up, so each tube is \emph{longer}: the time in every tube stays half an AC
        cycle. Limit: length (SLAC: \qty{3.2}{km}). Use: X-rays for cancer.
\end{itemize}

\subsection{The cyclotron}

\begin{sidebyside}[0.38]
  \centering
  \begin{tikzpicture}[x=0.88cm, y=0.88cm]
    \fill[hscseriesB!12] (0.15,-1.45) arc[start angle=-90, end angle=90, radius=1.75] -- cycle;
    \fill[hscseriesB!12] (-0.15,2.05) arc[start angle=90, end angle=270, radius=1.75] -- cycle;
    \draw[hscseriesB!75!black, line width=0.8pt] (0.15,-1.45) arc[start angle=-90, end angle=90, radius=1.75] -- cycle;
    \draw[hscseriesB!75!black, line width=0.8pt] (-0.15,2.05) arc[start angle=90, end angle=270, radius=1.75] -- cycle;
    \foreach \p in {(-1.45,1.25), (-1.45,-0.55), (-0.95,1.62), (-0.75,-1.25), (1.45,1.25),
                    (1.45,-0.65), (0.95,1.62), (0.75,-1.25), (-1.72,0.3)}
      \node[text=hscfield, font=\small] at \p {$\times$};
    \draw[hscvel, line width=0.9pt, postaction={decorate}, decoration={markings,
        mark=at position 0.12 with {\arrow{Stealth[length=2mm, width=1.6mm]}},
        mark=at position 0.45 with {\arrow{Stealth[length=2mm, width=1.6mm]}},
        mark=at position 0.80 with {\arrow{Stealth[length=2mm, width=1.6mm]}}}]
      (0,0) -- (0.15,0) arc[start angle=-90, end angle=90, radius=0.400] -- ++(-0.30,0)
      arc[start angle=90, end angle=270, radius=0.566] -- ++(0.30,0)
      arc[start angle=-90, end angle=90, radius=0.693] -- ++(-0.30,0)
      arc[start angle=90, end angle=270, radius=0.800] -- ++(0.30,0)
      arc[start angle=-90, end angle=90, radius=0.894] -- ++(-0.30,0)
      arc[start angle=90, end angle=270, radius=0.980] -- ++(0.30,0)
      arc[start angle=-90, end angle=90, radius=1.058] -- ++(-0.30,0)
      arc[start angle=90, end angle=270, radius=1.131] -- ++(0.30,0)
      arc[start angle=-90, end angle=90, radius=1.200] -- ++(-0.30,0)
      arc[start angle=90, end angle=270, radius=1.265] -- ++(0.30,0)
      arc[start angle=-90, end angle=-15, radius=1.327];
    \draw[velocity] (1.431,-0.009) -- (1.846,1.536);
    \node[font=\scriptsize, text=hscvel, anchor=west] at (1.86,1.5) {exit};
    \fill[hscink] (0,0) circle (1.6pt);
    \draw[hscink, line width=0.6pt] (-0.6,1.98) -- (-0.6,2.35) -- (-0.25,2.35);
    \draw[hscink, line width=0.6pt] (0.6,1.98) -- (0.6,2.35) -- (0.25,2.35);
    \draw[hscink, fill=white, line width=0.7pt] (0,2.35) circle (0.25);
    \node[font=\small] at (0,2.35) {$\sim$};
    \node[font=\scriptsize, anchor=west] at (0.62,2.42) {AC};
    \node[font=\scriptsize, text=hscseriesB!65!black] at (-1.15,-1.62) {dee 1};
    \node[font=\scriptsize, text=hscseriesB!65!black] at (1.15,-1.62) {dee 2};
    \node[font=\scriptsize] at (0,-1.62) {gap};
    \node[font=\scriptsize, align=center, anchor=north] at (0,-1.8)
      {$\times$: $B$ into the page \quad dot: ion source};
  \end{tikzpicture}
  \captionof{figure}{Top view of a cyclotron: a positive ion spirals out
    anticlockwise.}\label{ch5:fig:cyclotron}
\nextside
  A \term{cyclotron} has two hollow D-shaped \term{dees} in a uniform magnetic field, with
  AC across the gap (\cref{ch5:fig:cyclotron}).
  \begin{itemize}
    \item The gap's electric field speeds the ion (charged atom) up.
    \item In a dee (no electric field), $B$ bends it in a half circle.
    \item Faster means a bigger circle: the path spirals out.
    \item The AC flips while the ion is in a dee, so every crossing is a push.
    \item Limit: near $c$ each turn takes longer, so ions fall out of step. Uses: isotopes
          for PET (medical) scans; tumour therapy.
  \end{itemize}
\end{sidebyside}

\begin{formulacard}[notonsheet]{A charge moving in a circle in a magnetic field}
  \cardformula{r = \frac{mv}{qB} \qquad\qquad T = \frac{2\pi m}{qB}}
  \begin{symbols}
    \sym{$r$, $T$}{radius of the circle; period (time for one full circle)}{\si{m}, \si{s}}
    \sym{$m$, $v$}{mass and speed (speed at \ang{90} to $B$)}{\si{kg}, \si{m.s^{-1}}}
    \sym{$q$, $B$}{charge; magnetic field strength}{\si{C}, \si{T}}
  \end{symbols}
  \whentouse{well below $c$. $T$ does not depend on speed, so a cyclotron's AC has one
    fixed frequency, $f = 1/T$. Near $c$, use $r = p_\mathrm{v}/(qB)$.}
  \derive{magnetic force ($\theta = \ang{90}$) $=$ centripetal force (both on the sheet):
    $qvB = mv^2/r$, so $r = mv/(qB)$. Then $T = 2\pi r/v = 2\pi m/(qB)$ (the $v$
    cancels).}
  \trap{$r$ is the radius, not the diameter. Keep the $2\pi$ in $T$.}
\end{formulacard}

\subsection{The synchrotron and the Large Hadron Collider}

\begin{itemize}
  \item A cyclotron's magnet must cover the whole spiral. A \term{synchrotron} keeps the
        particles on one ring of fixed radius (after a linac and smaller rings).
  \item An \term{RF cavity} (a box with a fast-changing electric field) pushes them once per
        lap, in time (``synchro''). $B$ rises with the momentum, so $r = p_\mathrm{v}/(qB)$
        stays fixed.
  \item A \term{collider}: two beams meet head-on (total momentum zero), so all the kinetic
        energy can make new particles. \term{Large Hadron Collider} (LHC, CERN):
        \qty{27}{km} ring, \qty{8}{T} magnets, \qty{13.6}{TeV} (\qty{13.6e12}{eV}) per
        collision.
\end{itemize}

\subsection{What accelerators have shown}

A \term{detector} tracks what comes out: curved tracks give charge sign and momentum.
Masses are in \si{GeV}\chFiveperc\ (rest energy in \si{GeV}). Each row of
\cref{ch5:tab:tests} is a confirmed prediction; experiment can also lead theory (particle
zoo $\to$ quarks).

\begin{table}
  \caption{Standard Model predictions tested by accelerators.}\label{ch5:tab:tests}
  \centering
  \small
  \begin{tabular}{@{}L{55mm}L{50mm}L{35mm}@{}}
    \toprule
    \colhead{Theory predicts} & \colhead{Accelerator test} & \colhead{Result}\\
    \midrule
    the $\Omega^-$, its charge and mass (1962) & Brookhaven, 1964 & found, as predicted\\
    quarks inside protons (1964) & SLAC linac, 1968 & point-like charges\\
    W and Z, about 80 and \qty{91}{GeV}\chFiveperc & CERN collider, 1983 & found at those
      masses\\
    a third-generation quark (top) & Fermilab (USA), 1995 & found,
      \qty{173}{GeV}\chFiveperc\\
    the Higgs boson (1964) & LHC, CERN, 2012 & found, \qty{125}{GeV}\chFiveperc\\
    \bottomrule
  \end{tabular}
\end{table}

\subsection{Why our understanding is still incomplete}

\begin{itemize}
  \item \textbf{Gravity} is not in the Standard Model; the graviton has never been found.
  \item \textbf{Dark matter} (known only by its gravity) and dark energy (the unknown cause
        of the Universe's speeding-up expansion) are about 95\% of the Universe. No Standard Model particle explains them.
  \item \textbf{Antimatter:} the Big Bang should have made equal amounts of matter and
        antimatter, but the Universe is almost all matter.
  \item \textbf{Neutrino mass:} the first Standard Model had none; since 1998, a tiny mass.
  \item \textbf{Unknowns:} the reason for three generations; about 19 numbers must be measured.
\end{itemize}

%% ==================================================================
\section{Key experiments}\label{ch5:sec:experiments}
\covers{M8.IQ5.01, M8.IQ5.03}

\begin{experiment}{Deep inelastic scattering at SLAC (1967--1973)}
\textbf{Aim:} to find out whether the proton has parts inside it. \textbf{What was done:}
SLAC's \qty{3.2}{km} linac fired electrons (up to \qty{20}{GeV}) at liquid hydrogen. Magnetic
spectrometers (a magnet that bends the electrons, and detectors) measured each scattered
electron's energy and angle. \textbf{Result:} far more large-angle scattering than a smooth
proton allows. \textbf{Conclusion and significance:} protons and neutrons contain point-like
charged parts: quarks. This was the first direct evidence of quarks (Nobel Prize 1990).
\end{experiment}

\begin{experiment}{Discovery of the Higgs boson at the LHC (2012)}
\textbf{Aim:} to find the Higgs boson predicted in 1964, and measure its mass.
\textbf{What was done:} protons collided at \qtyrange{7}{8}{TeV} in total. The ATLAS and CMS
detectors kept the collisions that gave two photons or four leptons. From these, they found
the mass of the particle that made them. \textbf{Result:} a small bump at
\qty{125}{GeV}\chFiveperc\ above the background, in both detectors. The chance of a random bump
was below 1 in 3.5 million (``5 sigma''). \textbf{Conclusion and significance:} a new boson of
about \qty{125}{GeV}\chFiveperc\ exists: the Higgs boson. It was the last missing Standard Model
particle (Nobel Prize 2013 for Englert and Higgs, who predicted it).
\end{experiment}

%% ==================================================================
\begin{summarybox}{Remember}
\begin{itemize}
  \item SLAC's large-angle scattering showed quarks; beta decay and the $\Omega^-$ support it.
  \item Positron, muon, pions, antiproton, neutrino: each had a new mass, charge or behaviour.
  \item Quarks u, c, t: $+\tfrac23 e$; d, s, b: $-\tfrac13 e$. Baryon = 3 quarks; meson =
        quark + antiquark.
  \item Carriers: gluon (strong), photon (EM), $\mathrm{W}^\pm$ and $\mathrm{Z}^0$ (weak).
        Gravity is not in the model.
  \item $W = qV$ speeds particles up; $B$ only bends them. Cyclotron: $T = 2\pi m/(qB)$.
  \item Accelerators confirmed the $\Omega^-$, quarks, W, Z, top quark and Higgs boson.
  \item Gravity, dark matter, antimatter and neutrino mass: the model is incomplete.
\end{itemize}
\end{summarybox}
